% Modelo base — Currículo Camaleão
% Use o preâmbulo abaixo SEM alterar a estrutura. Substitua apenas o conteúdo
% entre os marcadores de seção (%% =========================).
% Regras de formatação por trás deste preâmbulo: ver regras-ats.md.

\documentclass[10pt,a4paper]{article}

\usepackage[utf8]{inputenc}
\usepackage[T1]{fontenc}
\usepackage[brazil]{babel}
\usepackage{lmodern}
\usepackage[margin=1.8cm]{geometry}
\usepackage{enumitem}
\usepackage{titlesec}
\usepackage{hyperref}

\hypersetup{
    colorlinks=true,
    urlcolor=black
}

\pagestyle{empty}
\setlength{\parindent}{0pt}
\setlength{\parskip}{0.4em}
\setlist[itemize]{leftmargin=1.3em, itemsep=0.2em}

\titleformat{\section}
  {\large\bfseries}
  {}
  {0pt}
  {}
  []

\titlespacing*{\section}{0pt}{1em}{0.3em}

\begin{document}

% ESTRUTURA OBRIGATÓRIA: a do currículo "Geral" do Daniel
% (Daniel_Vieira_Geral_Engenheiro_IA_Automacao_n8n.pdf), adotada pelo Cauê em
% 09/10/2026. Todo currículo sai nesta ordem, com estes blocos e estas
% categorias; só o conteúdo muda. Não crie seção nova (nem "Projetos"):
% projeto entra como bullet da experiência em que foi feito, e freelance é
% uma experiência própria ("Desenvolvedor Autônomo / Freelancer").

% Linha de título: variação de AI Full Stack Engineer + 1 ou 2 keywords da vaga
% entre parênteses, ex.: "AI Full Stack Engineer (Agentes de LLM | Python)".
\textbf{\Large [NOME COMPLETO]} \\
[TÍTULO ADAPTADO À VAGA]

[Cidade] - [Estado], [País] \\
Telefone: [+55 XX XXXXX-XXXX] \\
Email: [email@email.com] \\
LinkedIn: https://www.linkedin.com/in/[perfil] \\
GitHub: https://github.com/[usuario]

% =========================
\section*{Resumo Profissional}

[Um parágrafo: título + especialidades com as keywords da vaga, stack
principal, resultados com número e as ferramentas de destaque]

% =========================
\section*{Experiência Profissional}

\textbf{[Empresa]} \\
Cargo: [Cargo] \\
Período: [MM/AAAA -- Presente] \\
Local: [Cidade, Estado]

\begin{itemize}
\item [Bullet CAR com keyword da vaga e resultado quantificado]
\item [Bullet CAR com keyword da vaga e resultado quantificado]
\item [Bullet CAR com keyword da vaga e resultado quantificado]
\end{itemize}

% --- replicar para cada experiência, da mais recente para a mais antiga ---

% =========================
\section*{Competências Técnicas}

\textbf{IA e Agentes:} [...] \\
\textbf{Automação e Integração:} [...] \\
\textbf{Linguagens:} [...] \\
\textbf{Full Stack:} [...] \\
\textbf{Banco de Dados:} [...] \\
\textbf{Cloud e DevOps:} [...] \\
\textbf{Testes:} [...] \\
\textbf{Ferramentas e Metodologias:} [...]

% =========================
\section*{Formação Acadêmica}

\textbf{[Instituição]} \\
Curso: [Nome do Curso] \\
Período: [AAAA -- AAAA] \\
Status: [Em andamento / Concluído]

% =========================
% Única exceção à estrutura: sem certificação na base, a seção inteira sai
% (decisão do Cauê, 09/10/2026).
\section*{Certificações}

[Nome da Certificação (Instituição)] \\
[Nome da Certificação (Instituição) -- MM/AAAA]

% =========================
\section*{Idiomas}

Português: Nativo \\
Inglês: [Nível]

\end{document}
